\section{Klasyfikacja ataków promptowych}
Ataki promptowe można sklasyfikować między innymi według miejsca wstrzyknięcia złośliwej instrukcji. Wyróżnić można ataki pośrednie (ang. indirect prompt injection), w których system przyjmuje dane wejściowe ze źródeł zewnętrznych, przykładowo z plików, stron internetowych czy wiadomości e-mail zawierających treść, która może wpłynąć na zachowanie modelu w sposób nieoczekiwany bez bezprośredniego udziału użytkownika~\cite{owasp_llm_top10}. Obok ataków pośrednich, istotną podgrupą zagrożeń są również bezpośrednie ataki promptowe (ang. direct prompting attacks). Występują w scenariuszu, w którym atakujący jest głównym użytkownikiem systemu i wchodzi w interakcję z modelem poprzez interfejs zapytań, by skłonić model do zignorowania wbudowanych reguł bezpieczeństwa~\cite{nist_taxonomy_2025}.
Podklasą ataków bezpośrednich są ataki typu jailbreak (tłum. dosł. ucieczka z więzienia). Pojęcia prompt injection oraz jailbreak w literaturze bywają stosowane zamiennie, choć różnią się zakresem~\cite{owasp_llm_top10}. Wstrzykiwanie polecenia polega na manipulacji danymi wejściowymi w celu zmiany logiki i zachowania modelu, ale nie musi to obejmować ominięcia środków obronnych~\cite{owasp_llm_top10}. Z kolei ataki typu jailbreak koncentrują się na umożliwieniu nieuprawnionego użycia poprzez obejście wbudowanych w modelu mechanizmów odmowy, co może ułatwić realizację kolejnych, bardziej złożonych ataków~\cite{correia2026systematicliteraturereviewllm}.
Ataki promptowe można również kategoryzować na podstawie dostępu do modelu. Wyróżnić w ten sposób można ataki białoskrzynkowe, w ramach których atakujący ma pełny dostęp do parametrów modelu, danych treningowych, wag, czy architektury~\cite{chowdhury2024breakingdefensescomparativesurvey}. Atakom tym podlegają przede wszystkim modele udostępniane do pobrania z ogólnodostępnych źródeł, przykładowo z platformy Hugging Face. Natomiast w przypadku ataków czarnoskrzynkowych nie ma możliwości dostępu do parametrów modelu, więc atak przeprowadzany jest z minimalną wiedzą o systemie lub jej całkowitym brakiem~\cite{nist_taxonomy_2025}. Tego typu atak oparty jest na interakcji z modelem i analizie zwracanych przez niego odpowiedzi, co stanowi najbardziej prawdopodobny scenariusz ataku w praktycznych zastosowaniach modeli językowych.
Innym kryterium klasyfikacji ataków promptowych może być przebieg interakcji z modelem. Wyodrębnić można konwersacje jednoturowe (ang. single-turn attacks), których celem jest stworzenie pojedynczego, samowystarczalnego polecenia wejściowego do modelu, które natychmiastowo przełamuje zabezpieczenia modelu. Wyróżnić można też interakcje wieloturowe (ang. multi-turn attacks), w których atakujący wykorzystuje stan konwersacji i okno kontekstowe modelu, by przy użyciu sekwencji zapytań, krok po kroku skłonić model do naruszenia reguł bezpieczeństwa.
W dalszej części pracy analizie poddano metody bezpośrednich ataków promptowych. Omówiono w niej metody oparte na udostępnionych parametrach modelu, jak i manualne oraz automatyczne techniki czarnoskrzynkowe, uwzględniając przy tym interakcję jednoturową oraz wieloturową.